\section{Open Questions for Follow-Up}

The following five questions remain genuinely open after this audit.
Each is grounded in a specific gap identified in \S5 (Critique) or in
one of the data-blocked rows of the C1--C8 claim ledger. Each has
concrete next steps executable under LUCID's free-endpoint /
no-author-contact protocol.

\subsection*{Q1. Does modern SEER (1973--2020) still prefer the paper's
two-stage MVK best-fit over higher-stage variants?}

\textbf{Basis.} Claim C6 (SEER colon model comparison) is DATA-BLOCKED in the
present audit because the Little-Wright generalized k+m stage Poisson-
likelihood fitter has never been publicly released and SEER microdata is
SEER*Stat-registration gated. The paper's headline model-comparison
claim (2-stage best; 5-stage significantly worse at $P<0.01$) has never
been independently reproduced. Modern SEER18 adds 18 years of follow-up
plus the screening-era.

\textbf{Next steps.} Register for a free SEER*Stat account; pull SEER18
colon-adenocarcinoma incidence by sex, single-year age, single-year
year, 1973--2020; re-implement the Little-Wright k+m stage Poisson-
likelihood ($\approx 200$ lines Python + \texttt{scipy.optimize});
fit $k\in\{2,3,4,5\}$, $m\in\{0,1,2\}$; compare log-likelihoods and
Akaike weights. Report whether the paper's ordering survives extended
follow-up + screening heterogeneity.

\subsection*{Q2. Is the SVM bystander-protection U-shape robust to the
competing-risk formulation choice?}

\textbf{Basis.} Both the paper (Talk 4) and REPORT.md claim C4 treat the
SVM cell-death removal as constant-rate. Whether the U-shape is a
genuine dose-response feature or an artifact of that constant-hazard
assumption has not been probed. The LUCID audit only recovered the
U-shape qualitatively at the paper's quoted $k_{\mathrm{ap}}$ values.

\textbf{Next steps.} In the existing \texttt{code/smoke\_replication.py}
SVM stub, replace the constant cell-death rate with a two-parameter
Gompertz form $h_{\text{death}}(t) = h_0 \exp(g t)$; sweep
$g\in\{0, 0.005, 0.02, 0.05\}$~d$^{-1}$; re-plot at
$D\in\{0,5,10,20,50,100,500\}$~mGy; measure U-shape dip depth vs $g$;
check whether the delayed-vs.-immediate $k_{\mathrm{ap}}$ ordering
(0.054 vs.\ 0.022 d$^{-1}$) survives. If the U-shape vanishes for
$g>0.02$~d$^{-1}$, that weakens the paper's protective-bystander
headline. Pure analytic; no external data.

\subsection*{Q3. Do the paper's MVK / TSCE hazard predictions agree with
modern low-dose radiation-protection benchmarks (BEIR-VII, UNSCEAR
2020)?}

\textbf{Basis.} The paper stops at qualitative claims about MVK hazard
shape. Neither the paper nor REPORT.md ever tabulates ERR/EAR
predictions against BEIR-VII Table 12D-1 or UNSCEAR 2020 Annex B. The
gap between mechanistic prediction and consensus empirical fit is
operationally central to low-dose radiation protection.

\textbf{Next steps.} Using \texttt{code/claim\_audit.py}'s MVK plateau
formula, compute predicted lifetime attributable risk (LAR) per Sievert
for solid cancers at age-at-exposure $\{10,30,50,70\}$ both sexes,
chronic 100~mSv exposure; compare to BEIR-VII Table 12D-1 and UNSCEAR
2020 Annex B site-specific tables (both free-download PDFs); report
percent deviation per age-at-exposure bin; flag mechanistic model as
insufficient if deviations exceed a factor of 2 at diagnostic-dose
ranges.

\subsection*{Q4. Does the paper's TSCE-implied Thorotrast liver-cancer
individual heterogeneity match modern polygenic-risk-score (PRS)
distributions from HCC GWAS?}

\textbf{Basis.} Claim C8 is DATA-BLOCKED because the Heidenreich-Luebeck-
Hazelton Thorotrast fitted posterior was never released. But modern
liver-cancer PRS distributions from UK Biobank HCC GWAS are published
as open-access summary statistics. Comparing the two heterogeneity
distributions is a novel cross-check unavailable in 2007.

\textbf{Next steps.} Pull UK Biobank HCC GWAS summary statistics (GWAS
Catalog, open access); compute population log-relative-risk
distribution under a standard PRS construction; compare interquartile
range and top-percentile ratio to the paper's ``$\geq 95\%$ at age 40
below $10\%$ of mean; top percentile $>10\times$ mean.'' If PRS spread
is markedly narrower than TSCE-implied heterogeneity, the TSCE fit is
absorbing unmodeled environmental variance into apparent genetic
heterogeneity.

\subsection*{Q5. Does an ML meta-model over multi-cohort radiation
epidemiology (LSS, INWORKS, Techa River, Mayak) outperform the paper's
MVK/TSCE family in cross-cohort risk prediction?}

\textbf{Basis.} The paper (2008) explicitly limits itself to mechanistic
biological models. INWORKS 2015 and 2023 updates plus Techa River and
Mayak follow-ups postdate the paper and would establish whether the
mechanistic constraint helps or hurts modern predictive accuracy.

\textbf{Next steps.} Assemble published summary-level ERR-per-Gy from
LSS (Preston 2007, Grant 2017), INWORKS (Richardson 2015, 2023), Techa
River (Krestinina), Mayak (Sokolnikov) --- all open-access tables; fit
leave-one-cohort-out cross-validated gradient-boosted regressor on
age-at-exposure, sex, time-since-exposure, dose; compare cross-cohort
predictive RMSE against MVK/TSCE closed-form predictions from
\texttt{code/smoke\_replication.py}; inspect SHAP values to see if the
ML model recovers MVK's two-stage age dependence or identifies unmodeled
interactions. No microdata required.
